\section{引言}
\label{sec:intro}

%
格点场论与统计力学中的一个关键问题是对场构型的积分求值，即所谓的路径积分。通常，这类积分通过马尔可夫链蒙特卡洛（MCMC）方法求值：利用一条马尔可夫链，从由理论作用量决定的期望概率分布中抽取场构型样本。
%
一个重要的实际问题是链中构型之间相关性的存在。临界慢化（critical slowing down）\cite{Wolff:1989wq} 指的是当逼近参数空间中的临界点时，相应的自相关时间发散。这种行为会急剧增加这些参数区域中模拟的计算代价~\cite{DelDebbio:2004xh,Meyer:2006ty}。对某些模型，人们已找到能显著减轻或消除这种慢化的算法~\cite{Kandel:1988zz,Bonati:2017woi,PhysRevD.98.054502,SwendsenWang87,Wolff1989,ProkofevSvistunov01,KawashimaHarada04,Bietenholz1995MeronCluster}，从而实现高效模拟。对场论而言，人们提出了多种借助杂交蒙特卡洛（Hybrid Monte Carlo, HMC）技术的各种变体~\cite{Ramos:2012bb,Gambhir:2015rha,Cossu:2017eys,Jin2019QNHMC}、多尺度更新流程~\cite{Endres:2015yca,Detmold:2016rnh,Detmold:2018zgk}、开边界条件或不可定向流形~\cite{Luscher:2012av,Mages:2015scv,Burnier:2017osu}、元动力学（metadynamics）~\cite{Laio:2015era} 以及机器学习工具~\cite{Tanaka:2017niz,Shanahan:2018vcv} 来绕过临界慢化的方案。然而在一些重要的理论类别中，临界慢化仍是限制性因素；例如在量子色动力学的格点表述（QCD，标准模型中描述强相互作用的部分）里，它是实现精确控制连续极限所需的细格距模拟的一大障碍。

本文提出一种新的\emph{基于流的 MCMC} 方法并将其应用于格点场组态的生成。所得马尔可夫链的自关联性质可以通过采样之前的一个优化（训练）步骤系统地加以改善。
在该方法中，先从一个简单分布抽取样本 $z$，再经变量替换（或称``流''）$\phi = f^{-1}(z)$ 变换，得到具有新有效分布 $\netp$ 的样本 $\phi$。映射 $f^{-1}$ 的选取要求计算高效，使样本 $\phi$ 易于抽取；并在一个变分族内优化，使分布 $\netp$ 逼近期望分布。
为保证采样的渐近精确性，以 $\tilde{p}_{f}$ 作为提议分布构造了带 Metropolis-Hastings 步骤的马尔可夫链。由于提议样本与链中先前样本相互独立，自相关时间与接受率相互耦合；当接受率趋于 $1$ 时自相关时间降为零。即便在标准算法呈现临界慢化的参数区域也是如此。在温和条件下（详见第~\ref{sec:flowMCMC} 节），该方法被保证在大量更新的极限下产生来自期望概率分布的样本。

该方法具有若干使其适合于格点场论路径积分求值的特点：
%
\begin{enumerate}
    \item 通过训练模型可系统地降低马尔可夫链的自相关时间；
    \item
    马尔可夫链的每一步只需一次模型求值和一次作用量计算；
    \item 每个更新提议都与前一个样本无关，因此提议可以并行生成并高效地串接成马尔可夫链；
    \item 模型使用模型自身产生的样本进行训练，无需来自期望概率分布的现成样本。
\end{enumerate}

其他一些机器学习方法也已被应用于 MCMC，对象包括统计力学系统、合成分布以及简单的格点量子场论。自学习蒙特卡洛（Self-learning Monte Carlo, SLMC）方法已被较为成功地应用于一到三维的 Ising 系统与费米子系统。这些方法通过多种技术为理论构造一个比原哈密顿量更容易抽样的有效哈密顿量~\cite{Huang2017,Liu2017SLMC,Liu2016SLMCFermion, Nagai2017SLMCCT,Shen2018SLMCDNN}。
有效哈密顿量的学习采用监督学习技术，其训练数据来自现有 MCMC 模拟、随机变异样本以及加速后的马尔可夫链本身的组合（因此得名``自学习''）。
基于流的方法已被用于二维 Ising 模型~\cite{Li2018}、多体系统~\cite{NoeBoltzmannGenerators2018} 以及
合成分布~\cite{Song2017,Levy2017GenHMC} 的蒙特卡洛抽样；生成对抗方法则被应用于二维标量场论~\cite{Urban2018,Zhou2018}。

与这些方法不同，本文提出的方法聚焦于直接从真实分布的一个良好近似生成样本，且每个生成样本的精确似然均已知。直接生成使得可以像 SLMC 那样自学习；而已知似然则允许使用 Metropolis-Hastings 接受步骤来保证精确性。


本文第~\ref{sec:flowMCMC} 节详述所提出的基于流的 MCMC 算法。第~\ref{sec:phi4} 节给出其在二维 $\phi^4$ 理论中的有效性数值研究。最后，第~\ref{sec:summary} 节概述将该方案推广到更多时空维数定义的理论以及更复杂场论（如 QCD）所需的发展与规模扩展。
